% R. Brown's translations of two letters to L. Breen, with his note and comments (cote 134-2): archive cote 134-2, site page 134 (PDF page 135).
% translation of the letter of 17-19 July 1975. Draft compared: ps1.tex, lines 1563-1572.
% AI-transcribed by Claude Opus 5.5 (claude-opus-5-5, adaptive thinking, effort high), Message Batches API, batch brown_full (msgbatch_01Md3i5yK1PxkpJn5GHfEExE), 2026-11-13; not yet checked by eye.
\scanpage{134-2/134}
\mspage{- 7 -}

I nevertheless understand your instinctive resistance to
conceive this extreme stripping of a beautiful topos $X$ , to the
point of retaining only the meagre homotopy type. Even more I
am persuaded that going to the root of this instinctive resistance,
one arrives at a generalisation and deepening of the notion of
``homotopy type'', which bring us to still new waters in the
development of a good homotopical yoga. Here is what I wish to
come to pass.

Let us speak first of sheaves (of sets, or of modules, etc.)
instead of stacks, for simplicity, and place ourselves in the
étale topos of a scheme. The locally constant sheaves - modulo
a supplementary condition of finiteness which is sufficiently
anodyne - form the easiest of the \uline{constructible} sheaves, for
the definition of which they serve as models. Supposing $X$
quasi-compact and quasi-separated, the general constructible
sheaves are those for which there exist a finite partition
$X = \hand{\bigcup}_{i\in I} X_i$ of $X$ in ``cells'' or ``strata'' $X_i$ , each locally
closed and constructible, such that the restriction of $F$ to
every $X_i$ is locally constant (also a finiteness condition ...).
Thus the category of constructible sheaves on $X$ (which gives
back the category of all sheaves on passing to a category of
ind-objects ...) appears itself to be an inductive limit of
categories associated to finer and finer partitions on $X$ . One
can then, by such a finite partition $P$ , propose a more profound
study of the category of sheaves (or complexes of sheaves, or
stacks) which are ``P-constructible'' (or, more generally, which are

\mspage{\hand{135}}
